\section{Introduction}
\label{sec:introduction}

Generative artificial intelligence is widely expected to be the defining labor-market shock of the 2020s. The launch of ChatGPT on 30 November 2022 gave the world's first general-purpose large language model (LLM) to hundreds of millions of users within months, and occupational exposure measures suggest that roughly half of US jobs have at least one task that LLMs can perform twice as fast as a human \citep{Eloundou2023_gpt_exposure}. Against these projections, the first set of careful causal estimates is striking. \citet{HumlumVestergaard2025_llm} use Danish administrative data linked to a 100{,}000-worker adoption survey and rule out wage and hours effects larger than two percent, two years after launch. Two-percent effects are economically small. They are also geographically narrow: essentially every credible causal estimate of how LLMs have reshaped aggregate labor markets comes from the United States, Denmark, or online freelance platforms in which workers in advanced economies dominate.

This leaves the first-order question - whether the most widely hyped general-purpose technology of the decade has actually moved labor-market outcomes - essentially unresolved for the developing world. Latin America is the natural place to ask. The region's aggregate informality rate sits near 55 percent, with country-specific rates between 25 percent in Uruguay and more than 70 percent in parts of the Andean region. Formal and informal workers occupy segmented markets with different wage-setting institutions, different degrees of exposure to digital technology, and different buffers against productivity shocks. Exposure to LLMs in Latin America is also structurally different: \citet{Azuara2024_idb_ai} show that predicted exposure tilts sharply toward formal, female, urban, and tertiary-educated workers in Chile, Mexico, and Peru, and \citet{Ciaschi2025_distributive_ai} and \citet{EganadelSol2025_ai_lac} extend this exposure-only picture to the broader region. What no existing paper provides is a causal estimate of whether this exposure has translated into actual labor-market responses.

This paper fills that gap. I assemble a harmonized panel of seven Latin American labor force surveys (Peru, Uruguay, Mexico, Chile, Colombia, Costa Rica, and Ecuador) covering 2021Q1 through 2025Q4, map Eloundou et al.'s occupational exposure scores from SOC-2010 to ISCO-08 via the BLS crosswalk (match rate 98.2 percent at the 4-digit level), and estimate the Average Causal Response on the Treated (ACRT) along the continuous exposure dimension using the difference-in-differences estimator of \citet{CallawayCGBS2024_continuous}. The post-COVID time window is a deliberate design choice: including pre-2020 periods would pool labor markets in two structurally different equilibria, violating the comparability requirement for the DiD control group. ChatGPT's launch delivers a single adoption date common to all occupations; identification comes from within-(occupation, country) variation in how exposure interacts with the post-launch indicator. Baseline (2021) formality status is pre-registered as the paper's central heterogeneity margin, on the argument that contemporaneous formality is itself an outcome of LLM exposure and therefore a bad control.

% RESULTS TBD — four narratives pre-registered, one realized
The results fall into one of four pre-committed narratives. I state the narratives now and defer the quantitative headline numbers to \cref{sec:results}, where they appear with exposure-normalized effect sizes, confidence intervals, and Honest-DiD bounds. \emph{Narrative A} (informality amplifies): aggregate LAC effects exceed two percent and are driven by the informal channel, with $\tau_I < \tau_F < 0$. \emph{Narrative B} (LAC-Denmark null replication): aggregate effects are indistinguishable from zero in both channels. \emph{Narrative C} (null in formal, large in informal): $\tau_F \approx 0$ and $\tau_I$ is meaningfully negative, identifying a channel that Danish and US designs cannot detect because their informal sectors are negligible. \emph{Narrative D} (stratified gains): positive wage effects concentrate in tertiary-educated urban formal workers, bypassing the informal majority.

% Placeholder finding paragraphs to be activated once results arrive.
% [Narrative A finalization paragraph]: Consistent with Narrative [X]. On the full sample, a one-standard-deviation increase in Eloundou ζ exposure is associated with a [X] pp change in employment and a [X] percent change in log hourly wages post-2023. The informal channel accounts for [X] percent of the aggregate effect.

The paper's cleanest contribution is its contrast with \citet{HumlumVestergaard2025_llm}. Denmark's combination of low informality, high digital readiness, strong worker bargaining, and generous social insurance is precisely the structural configuration most likely to mute a labor-market response: buffers absorb the shock before it registers in administrative earnings data. Latin America's configuration is the opposite - high informality, fragmented wage-setting, weaker social insurance, and lower but non-trivial adoption. If the Danish null is universal, it should survive in this very different setting; if it reflects Danish institutions rather than LLM capabilities, the LAC estimates should diverge in a predictable direction. The seven-country panel sharpens this comparison by delivering within-region variation in informality (25 to 70 percent), digital penetration, and ChatGPT adoption intensity.

The mechanism is organized around the task-based framework of \citet{AcemogluRestrepo2018_race} and \citet{AcemogluAutor2011_tasks}. LLMs act simultaneously through a displacement channel (substituting for non-routine cognitive tasks) and a reinstatement channel (complementing human judgment in tasks that require contextual reasoning). The net effect on employment and wages is empirical, not theoretical, and its decomposition across formality strata is exactly what the pre-registered specification targets. LLMs also break the classical routine-biased technical change pattern of \citet{AutorLM2003_skill_content}: the tasks most affected (writing, synthesis, analysis) are non-routine cognitive tasks historically considered protected from automation.

The paper makes five contributions. First, it provides the first causal estimates of LLM exposure on labor-market outcomes for Latin America, filling the gap between the descriptive exposure literature \citep{Azuara2024_idb_ai, Ciaschi2025_distributive_ai, EganadelSol2025_ai_lac, WorldBank2024_genai_lac} and the OECD-focused causal literature \citep{HumlumVestergaard2025_llm, Hartley2024_labor_effects, Chen2025_llm_unemployment}. Second, it applies \citet{CallawayCGBS2024_continuous} continuous-treatment DiD to LLM exposure; the prior causal literature has relied on binary cutoffs, discarding most of the dose variation. Third, it elevates baseline formality from a subgroup check to the co-primary identification margin, with an imputation protocol that avoids using a post-treatment outcome as a conditioning variable. Fourth, it builds a reproducible seven-country harmonized crosswalk and panel (2021--2025) that future LAC researchers can extend. Fifth, the paper pre-commits to four narratives before the analysis is run, so that the narrative emphasized in the final draft is determined by which empirical pattern obtains rather than by post-hoc specification search.

The rest of the paper proceeds as follows. \Cref{sec:literature} places the contribution in relation to the AI exposure literature, the post-ChatGPT causal evidence, the task-based theoretical framework, and the LAC-specific context. \Cref{sec:method} presents the identification strategy, the continuous-DiD estimating equation, the four pre-registered narratives, and the robustness layer. \Cref{sec:data} describes the seven-country harmonized panel. \Cref{sec:results} reports the main estimates and the baseline-formality decomposition. \Cref{sec:robustness} presents the robustness evidence, and \Cref{sec:conclusion} discusses external validity and policy implications.
